\section{Splitting at a scale}
\label{sec:splitting}

The first operation below splits the graph into small boxes, allowing a
small amount of vertex overlap.  The second operation extracts expanding
pieces on disjoint vertex sets within each box.  Only edges leave when a
piece is extracted: the residual graph retains every box vertex.  Thus
every earlier piece vertex still occurs in a descendant box when the
operations are iterated.  The overlapping split
and its inverse-logarithmic order potential adapt the expander
splitting argument of Buci\'c and Montgomery~\cite[Lemma~14]{BM2024}.
Here we stop at a prescribed scale and obtain a total box order
of \(N(1+O(1/\log D))\), which will be charged in the final induction.

\begin{lemma}[Splitting at a scale]\label{lem:scale-split}
There are absolute constants $D_{\mathrm{sp}}$ and $C_0$ such that, for
every real $D\ge D_{\mathrm{sp}}$, every graph $G$ on a nonempty vertex
set of order $N$ and with average degree at most $D$ admits the following
two successive operations.  One may take $C_0=2$.
\begin{enumerate}[label=\textup{(\roman*)}]
\item\label{item:scale-boxes}
Remove a family $\mathcal C$ of edge-disjoint cycles, each of length at
least $D^{1.01}$, where
\begin{equation}\label{eq:scale-cycle-cost}
  \abs{\mathcal C}\le \frac{N}{2D^{0.01}}.
\end{equation}
Partition the remaining edges into graphs $G_Y$ on an indexed family
$\mathcal Y$ of vertex sets, called boxes.  These satisfy
\begin{equation}\label{eq:scale-box-order}
\begin{gathered}
  \bigcup_{Y\in\mathcal Y}Y=V(G),\qquad
  1\le \abs Y\le D^{1.02},\\
  \sum_{Y\in\mathcal Y}\abs Y
    \le \left(1+\frac{C_0}{\log D}\right)N.
\end{gathered}
\end{equation}
The graphs $G_Y$ contain only their assigned edges; their vertex sets
may overlap.

\item\label{item:scale-expanders}
Put $\sigma=D^{0.90}$.  For each box $Y$, partition $E(G_Y)$ into
the edge sets of a family $\mathcal R_Y$ of graphs $R$ and a graph $H_Y$.
The vertex sets of the graphs in $\mathcal R_Y$ are pairwise disjoint
subsets of $Y$.  Each $R$ has cut expansion at least $\sigma$ and order
\begin{equation}\label{eq:scale-piece-order}
  \sigma\le \abs{V(R)}\le D^{1.02}.
\end{equation}
The residual graph $H_Y$ is taken on all of $Y$, including isolated
vertices, and satisfies
\begin{equation}\label{eq:scale-residual-degree}
  \frac{2\abs{E(H_Y)}}{\abs Y}\le D^{0.95}.
\end{equation}
\end{enumerate}
Thus the cycles in $\mathcal C$, the pieces $R$, and the residual graphs
$H_Y$ together partition $E(G)$.
\end{lemma}

We first give the elementary long-cycle argument used to find the
overlapping splits. It is the depth-first-search and bypass argument
of Buci\'c and Montgomery~\cite[Lemma~25]{BM2024}, with constants
suited to the present scale cutoff. Related depth-first-search balance
arguments appear in Ben-Eliezer, Krivelevich and
Sudakov~\cite[Lemma~2.3]{BKS2012}; Krivelevich~\cite[Theorem~1]{Krivelevich2019}
gives a stronger general criterion for long cycles from local expansion.
For a vertex set $U$ in a graph $H$, write
$N_H(U)=\bigcup_{v\in U}N_H(v)$; its external neighbor set is
$N_H(U)\setminus U$.

\begin{lemma}[Vertex expansion gives a long cycle]
\label{lem:split-vertex-expansion}
Fix $\eps=1/32$.  For every sufficiently large integer $u$, put
$\alpha=\eps/\log^2 u$.  If a graph $H$ of order $u$ satisfies
\begin{equation}\label{eq:split-vertex-expansion}
  \abs{N_H(U)\setminus U}\ge\alpha\abs U
  \qquad\left(0<\abs U\le\frac u2\right),
\end{equation}
then $H$ contains a simple cycle of length greater than
\begin{equation}\label{eq:split-long-cycle-length}
  q(u)=\frac{\alpha^2u}{100}
      =\frac{\eps^2u}{100\log^4u}.
\end{equation}
\end{lemma}

\begin{proof}
The expansion condition implies that $H$ is connected: in a
disconnected graph some component has order between $1$ and $u/2$ and
has no external neighbors.

Run a depth first search.  Its state consists of the unexplored
vertices, a stack whose vertices in stack order form a simple path, and the finished
vertices.  Initially the stack consists of one vertex, all other
vertices are unexplored, and none are finished.  If the top vertex of
the stack has an unexplored neighbor, push such a neighbor onto the
stack.  Otherwise pop the top vertex and mark it finished.  At all
times there are no edges between finished and unexplored vertices:
when a vertex is finished it has no unexplored neighbor, and the
unexplored set subsequently only shrinks.  Connectivity ensures that
the stack cannot become empty while an unexplored vertex remains, so
this search finishes all vertices.

The number of unexplored vertices minus the number of finished
vertices starts at $u-1$, ends at $-u$, and decreases by exactly one at
every push or pop.  Consequently there is a state at which these two
numbers are exactly equal, say to $t$.  Let $P$ be the stack path at
that state, and let $L=\abs{V(P)}=u-2t$.  If $t\ge u/3$, then
$0<t\le u/2$, and all external neighbors of the unexplored set lie on
$P$.  Applying \eqref{eq:split-vertex-expansion} gives
$L\ge\alpha t\ge\alpha u/3$.  If $t<u/3$, then
$L>u/3\ge\alpha u/3$, since $\alpha\le1$ for the values of $u$ under
consideration.  Thus in either case
\begin{equation}\label{eq:split-stack-order}
  L\ge\frac{\alpha u}{3}.
\end{equation}

Write $q=q(u)$ and $j=\floor q$.  We may assume $q\ge1$, since
$u/\log^4u$ tends to infinity.  In particular $j\ge1$, and
\[
  j+1\le2q=\frac{\alpha^2u}{50}
       \le\frac{\alpha u}{6}\le\frac L2.
\]
Partition $P$ into three consecutive vertex segments $X,J,Z$, with
$\abs J=j$, distributing the other $L-j$ vertices as evenly as
possible between $X$ and $Z$.  Both outer segments have order at least
\[
  \floor{\frac{L-j}{2}}
    \ge\frac{L-j-1}{2}\ge\frac L4.
\]
These inequalities include the integer rounding needed to make the
three segments nonempty.

Suppose first that $H-J$ has a path from $X$ to $Z$.  Choose one of
minimum length among all such paths, with endpoints $x\in X$ and
$z\in Z$.  A shortest path is simple.  Its interior contains no vertex
of $X$ or $Z$, because encountering either set internally would give a
shorter path between the two sets.  It also avoids $J$.  Since
$X,J,Z$ partition $V(P)$, the chosen path meets $P$ only at $x,z$.
Its union with the $x$--$z$ subpath of $P$ is therefore a simple cycle.
That subpath contains all $j$ vertices of $J$ and has at least $j+1$
edges; the new path has at least one edge.  The cycle consequently has
at least $j+2\ge3$ edges, and its length is greater than $q$.

It remains to show that such a path must exist.  Otherwise $X$ and
$Z$, each connected by its segment of $P$, lie in distinct components
of $H-J$.  Separate the component containing $X$ from the union of all
other components, which contains $Z$, and let $A$ be the smaller of
these two sets.  Then
\[
  \frac L4\le\abs A\le\frac{u-j}{2}\le\frac u2,
  \qquad N_H(A)\setminus A\subseteq J.
\]
Expansion and \eqref{eq:split-stack-order} imply
\[
  j\ge\alpha\abs A\ge\frac{\alpha L}{4}
    \ge\frac{\alpha^2u}{12}
    >\frac{\alpha^2u}{100}=q,
\]
contrary to $j=\floor q$.  This proves the lemma.
\end{proof}

\begin{proof}[Proof of \Cref{lem:scale-split}]
We prove the two operations in order.  All lower bounds imposed on $D$
in this proof are absolute and hold simultaneously for every
$D\ge D_{\mathrm{sp}}$ after one choice of $D_{\mathrm{sp}}$.

\paragraph{Removing long cycles.}
Choose a maximal family $\mathcal C$ of edge-disjoint cycles in $G$
whose lengths are at least $D^{1.01}$.  Since
$\abs{E(G)}\le ND/2$, we have
\[
  \abs{\mathcal C}D^{1.01}\le\abs{E(G)}\le\frac{ND}{2},
\]
which proves \eqref{eq:scale-cycle-cost}.  Integral cycle lengths are
compared with the real threshold as stated, equivalently with
$\ceil{D^{1.01}}$.  Delete these cycle edges and retain the entire
vertex set.  By maximality the remaining graph has no cycle of length
at least $D^{1.01}$.  Every graph subsequently formed in the first
operation is a subgraph of this remainder, so it inherits this
prohibition.

\paragraph{A sparse external neighborhood.}
Put $T=D^{1.02}$ and $\eps=1/32$.  We claim that every current graph
$H$ of order $u>T$ has a nonempty vertex set $U$ such that
\begin{equation}\label{eq:split-sparse-neighborhood}
  \abs U\le\frac u2,\qquad
  \abs{N_H(U)\setminus U}
    <\frac{\eps\abs U}{\log^2u}.
\end{equation}
Otherwise \Cref{lem:split-vertex-expansion} supplies a cycle longer
than $q(u)$ as defined in \eqref{eq:split-long-cycle-length}.
To obtain a contradiction uniformly for every $u>T$, including $u$
arbitrarily large relative to $D$, observe that $x/\log^4x$ is
increasing for $\ln x>4$, as follows by differentiating.  Moreover,
\begin{equation}\label{eq:split-uniform-margin}
  \frac{q(T)}{D^{1.01}}
    =\frac{\eps^2}{100(1.02)^4}
       \frac{D^{0.01}}{(\log D)^4}
    \longrightarrow\infty\qquad(D\longrightarrow\infty).
\end{equation}
Choose $D_{\mathrm{sp}}$ large enough that $T$ is in the range of
\Cref{lem:split-vertex-expansion}, that $T>e^4$, and that
$q(T)>D^{1.01}$.  Then $q(u)>D^{1.01}$ for every $u>T$,
contradicting the inherited cycle prohibition.  This proves
\eqref{eq:split-sparse-neighborhood}.

\paragraph{The overlapping split and its edge assignment.}
For such a set $U$, put $B=N_H(U)\setminus U$ and form two children
on vertex sets
\[
  V_1=U\cup B,\qquad V_2=V(H)\setminus U.
\]
Assign to the first child every edge of $H$ with both endpoints in
$V_1$, and assign to the second child every other edge of $H$.
Every edge incident with $U$ has both endpoints in $U\cup B$.
Consequently every edge assigned to the second child has both endpoints
in $V_2$.  This is an edge partition into two simple subgraphs; in
particular, edges with both endpoints in the overlap $B$ are assigned
only to the first child.  The child vertex sets cover the parent and
have intersection exactly $B$.  All vertices of each specified child
set are retained, even when isolated in its assigned graph.

Writing $s=\abs U$, $b=\abs B$, $u_1=s+b$ and $u_2=u-s$, we have
$1\le s\le u/2$ and $b<\eps s/\log^2u$.  For sufficiently large
$D$, therefore,
\begin{equation}\label{eq:split-child-orders}
\begin{gathered}
  1\le u_1\le\frac{3u}{4},\qquad 1\le u_2<u,\\
  u_1+u_2=u+b
      \le u+\frac{\eps u_1}{\log^2u}.
\end{gathered}
\end{equation}
Indeed $\eps/\log^2u\le1/2$ gives
$u_1<(3/2)s\le3u/4$, and $u_2<u$ follows from $s\ge1$.
Repeat this split whenever the current order exceeds $T$.
Both child orders are positive integers strictly smaller than their
parent order, so the resulting binary splitting tree is finite.
Its leaves give the graphs $G_Y$ on boxes of order at most $T$.
The child cover property gives $\bigcup_Y Y=V(G)$, and the edge
assignments at each split show that the leaf graphs partition exactly
the edges left after the cycle removal.

\paragraph{Controlling the sum of box orders.}
The sum of child orders can exceed the parent order by $b$.  The
following potential pays for this overlap.  For every real $x\ge1$
define
\begin{equation}\label{eq:split-potential}
  f(x)=1+\frac1{\log(T/2)}
         -\frac1{\log(\max\{x,T/2\})}.
\end{equation}
For $T\ge4$, this is nondecreasing and
\begin{equation}\label{eq:split-potential-range}
  1\le f(x)\le1+\frac1{\log(T/2)}\le2.
\end{equation}
At a split of a graph of order $u>T$, put
$M=\max\{u_1,T/2\}$.  By \eqref{eq:split-child-orders} and $u>T$,
we have $M\le3u/4$.  Both logarithms below are positive, and
\begin{equation}\label{eq:split-potential-decrement}
\begin{aligned}
  f(u)-f(u_1)
    &=\frac1{\log M}-\frac1{\log u}\\
    &=\frac{\log(u/M)}{\log M\,\log u}
      \ge\frac{\log(4/3)}{\log^2u}.
\end{aligned}
\end{equation}
Monotonicity also gives $f(u_2)\le f(u)$.  Hence, using
\eqref{eq:split-child-orders},
\begin{equation}\label{eq:split-potential-subadditivity}
\begin{aligned}
  u_1f(u_1)+u_2f(u_2)
    &\le (u_1+u_2)f(u)
         -\frac{u_1\log(4/3)}{\log^2u}\\
    &\le uf(u)
         +\frac{u_1}{\log^2u}
                 \bigl(\eps f(u)-\log(4/3)\bigr)\\
    &\le uf(u).
\end{aligned}
\end{equation}
For the last step, \eqref{eq:split-potential-range} gives
$\eps f(u)\le1/16<\log(4/3)$.  Repeatedly applying
\eqref{eq:split-potential-subadditivity} in the finite splitting tree,
and using $f\ge1$ on its leaves, proves
\[
  \sum_{Y\in\mathcal Y}\abs Y
    \le\sum_{Y\in\mathcal Y}\abs Y f(\abs Y)
    \le Nf(N)
    \le N\left(1+\frac1{\log(T/2)}\right).
\]
This argument also applies if $N\le T$, when the tree has just its
root.  Since $T=D^{1.02}$ and $D\ge4$,
\[
  \frac1{\log(T/2)}
    =\frac1{1.02\log D-1}\le\frac2{\log D}.
\]
This proves \eqref{eq:scale-box-order} with $C_0=2$ and finishes
operation~\ref{item:scale-boxes}.

\paragraph{Extracting pieces of cut expansion $\sigma$.}
This repeated removal of sparse cuts follows the small-side charging
method used by Conlon, Fox and Sudakov~\cite[Lemma~3.1]{CFS2014}.
We give the accounting for the precise cut expansion and residual
average degree required here.
Fix a box $Y$, write $y=\abs Y$, and start with its assigned graph
$G_Y$.  Whenever a current graph on a vertex set $A$ of order at least
two fails to have cut expansion $\sigma=D^{0.90}$, choose a partition
$A=A_1\mathbin{\dot\cup}A_2$ into nonempty sets such that its cut has
fewer than $\sigma\min\{\abs{A_1},\abs{A_2}\}$ edges.  Put those
cut edges into the residual graph $H_Y$, and continue separately with
the edges internal to each of $A_1,A_2$.  Both child orders strictly
decrease, so this procedure terminates.  Its terminal vertex sets
partition $Y$.  The nonsingleton terminal graphs form
$\mathcal R_Y$; by the stopping rule each has cut expansion at least
$\sigma$.  If such a graph has order $r$, its singleton cuts show that
its minimum degree is at least $\sigma$.  Simplicity then gives
$r-1\ge\sigma$, and in particular \eqref{eq:scale-piece-order}.
Terminal singletons carry no edges and are omitted from
$\mathcal R_Y$.

To bound the edges put into $H_Y$, at each split charge their number
equally to the vertices in a smaller side, choosing either side in
case of equality.  Each charged vertex receives less than $\sigma$
at that split.  Whenever a given vertex is charged, the current set
containing it decreases in size by a factor of at least two; any
intervening splits can only decrease its size further.  If the vertex
is charged $k$ times, after its last charge its set has order at most
$y/2^k$ and at least one.  Thus $k\le\log y$.

All residual cut edges are distinct, because each subsequent split
uses only edges internal to a previous child.  Summing the charges
therefore gives
\begin{equation}\label{eq:split-cut-charge}
  \abs{E(H_Y)}\le\sigma y\log y.
\end{equation}
Every edge not set aside ends in exactly one terminal graph, and
singleton terminal graphs contain no edges.  Thus $H_Y$ and the
graphs in $\mathcal R_Y$ partition $E(G_Y)$ exactly.  Define
$V(H_Y)=Y$, retaining also the vertices of all terminal graphs.
For $y=1$ the residual has no edges.  For every $1\le y\le D^{1.02}$,
\eqref{eq:split-cut-charge} yields
\[
  \frac{2\abs{E(H_Y)}}{y}
    \le2\sigma\log y
    \le2.04D^{0.90}\log D
    \le D^{0.95},
\]
where the last inequality holds for all sufficiently large $D$.
This proves operation~\ref{item:scale-expanders} and the lemma.
\end{proof}

In particular, when the two operations are repeated on the residual
graphs, the boxes descending from any fixed box continue to cover
every vertex of that box.  This follows from the cover property of
each overlapping split and the fact that $H_Y$ always has vertex set
$Y$, irrespective of which incident edges were removed as cycles or
expanding pieces.
